
\subsubsection{6.1 Key Findings}

\textbf{The confounding reversal is the primary finding.} The transformation from p = 0.583 to p = 0.0053 on the same datasets, same outcome variable, and same statistical test — achieved by adding propensity matching — demonstrates that suppressor confounding operates in this data. High-FAIR datasets are systematically newer; newer datasets have had less time to accumulate runs; this creates a negative confound that suppresses the raw FAIR-reuse correlation. Matching on creation-year proxy, task type, and dataset size removes this suppressor. Prior literature reporting null or weak correlations between FAIR compliance and dataset adoption may reflect the same confounding mechanism rather than true null effects, though this inference requires verification on independent data.

\textbf{Findable sub-criteria specificity.} Aggregate FAIR threshold (p = 0.697) versus Findable IV (p = 0.0053) from identical matched pairs suggests that different FAIR dimensions have heterogeneous effects on discovery speed. Aggregate scores aggregate these heterogeneous effects, diluting the Findable signal. The implication — if confirmed at production scale — is that Findable-specific investments (persistent identifiers, metadata richness, search indexing) may be more consequential for discovery speed than generic FAIR compliance improvements.

\textbf{The 44-day median TTFR reduction.} For ML research with iteration cycles of weeks to months, a 44-day reduction in median time-to-first-use could translate to earlier citation opportunities and faster dataset quality feedback. However, this estimate derives from a synthetic cohort with only 35 matched pairs and a wide Cox CI [1.032, 9.672]; the true magnitude of any production-scale effect is currently unknown.

\subsubsection{6.2 Limitations}

\textbf{L1: Synthetic proof-of-concept cohort (n = 200, 35 matched pairs).} All survival analysis results are preliminary. The Cox CI [1.032, 9.672] spans nearly an order of magnitude. The minimum recommended sample for reliable survival analysis is substantially larger than 35 events per group. Production replication with real OpenML upload dates is the immediate prerequisite.

\textbf{L2: Proxy FAIR scores, not true F-UJI sub-criteria.} The proxy captures only structural/computational characteristics. F-UJI sub-criteria for Interoperability (schema compliance, community standards) and Reusability (license clarity, citation guidance) have no proxy equivalent in the quality metrics used. All quantitative estimates (HR = 3.159, p = 0.0053) reflect proxy-quality-to-TTFR associations, not true F-UJI Findable sub-criteria associations. The direction and significance may not replicate when true F-UJI scores are used.

\textbf{L3: Upload_date unavailability.} The OpenML bulk \texttt{list_datasets()} API does not return \texttt{upload_date}. This prevented: (a) enforcement of the post-2018 cohort scope restriction, (b) verification of assumption A1 (non-retroactive FAIR tagging), and (c) a production-scale mechanism analysis for H-M2. Individual dataset API calls can fetch upload_date at approximately 1 request/second (requiring roughly 83 minutes for 5,000 datasets) and constitute the primary prerequisite for production replication.

\textbf{L4: H-M3 and H-M4 not executed.} The Reusable-dimension dominance hypothesis (H-M3) and HuggingFace cross-repository analysis (H-M4) were not executed. Predictions P2 (Reusable β > 0.15, dominant) and P3 (HuggingFace documentation completeness r > 0.15 with downloads) are inconclusive. The multi-repository and multi-dimension scope of the original hypothesis is not addressable from the current evidence.

\textbf{L5: Proportional hazards violation.} The Schoenfeld test flags a potential PH assumption violation in the smoke-test cohort. The Cox HR = 3.159 may not be constant across the observation window; time-varying hazard models are warranted at production scale.

\textbf{L6: Relaxed caliper in smoke test.} A caliper of 0.8 SD (vs. the production target of 0.2 SD) was used to obtain sufficient matched pairs from the n = 200 cohort. Wider calipers permit more heterogeneous matches, potentially amplifying within-pair differences artificially. Results under the production caliper may differ.

\textbf{L7: Assumption A1 (non-retroactive FAIR tagging) unverified.} Without upload dates, the temporal relationship between FAIR scoring and experimental runs cannot be confirmed. If high-reuse datasets received retroactive FAIR tagging, the causal ordering would be reversed — FAIR compliance would reflect achieved reuse rather than predict it.

\textbf{L8: Assumption A2 (run counts as deliberate engagement) unverified.} Automated pipeline runs, course projects, or scraping could inflate OpenML run counts. No user-ID filtering or algorithm-variety check was performed.

\subsubsection{6.3 Alternative Explanations for the Confounding Reversal}

Three competing explanations for the p = 0.583 → p = 0.0053 reversal merit consideration:

\begin{enumerate}
\item \textbf{Genuine suppressor confounding} (plausibility: high): Low-FAIR datasets are older and have had more time to accumulate runs; high-FAIR datasets are newer. Matching removes the age advantage of low-FAIR datasets, revealing genuine FAIR-driven faster discovery. This is the interpretation advanced in this paper.
\end{enumerate}

\begin{enumerate}
\item \textbf{Synthetic data artifact} (plausibility: medium): With 35 matched pairs from a synthetic n = 200 cohort, significance may reflect the cohort generation parameters rather than any real phenomenon. This explanation cannot be ruled out without production replication.
\end{enumerate}

\begin{enumerate}
\item \textbf{Caliper relaxation artifact} (plausibility: medium): The smoke-test caliper (0.8 SD) permits wider matches than the production standard (0.2 SD), which may artificially amplify within-matched-pair differences.
\end{enumerate}

The methodological contribution — that matched designs are necessary for credible FAIR-outcome studies — holds independently of whether explanation 1, 2, or 3 dominates, because the argument rests on the structural presence of confounders, not on any specific effect size.

\subsubsection{6.4 Broader Impact}

Establishing matched designs as a methodological standard for FAIR-outcome research improves evidence quality for infrastructure investment decisions. If confirmed at production scale, the sub-criteria specificity finding provides actionable guidance that could be applied by OpenML and HuggingFace administrators to prioritize Findable-specific improvements over generic FAIR compliance.

However, results should not be interpreted as evidence that only Findable improvements matter. We tested one mechanism (discovery speed) for one dimension. Sustained research engagement may depend critically on Reusable properties (licenses, provenance) that the present study does not address. The proxy FAIR score also correlates partially with dataset size; administrators should not conflate larger datasets receiving higher proxy scores with findability driving discovery.

